\documentclass[11pt,a4paper]{article}

%% PACHETE MATEMATICE
\usepackage{amsmath,amssymb,amsfonts}
\usepackage{amsthm}
\usepackage{mathpazo}
\usepackage{eulervm}
\usepackage{enumerate}
\usepackage{color}
\usepackage{graphicx}
\usepackage[all]{xy}
\usepackage{xcolor}
\usepackage{framed}
\usepackage[unicode]{hyperref}
\setcounter{MaxMatrixCols}{20}
\usepackage{multicol}
% \usepackage[romanian]{babel}


%% ASPECT
\linespread{1.05}
\usepackage[T1]{fontenc}
\usepackage[utf8x]{inputenc}
\usepackage[margin=2cm]{geometry}
% \usepackage[color]{showkeys}
\usepackage{anyfontsize}
\usepackage{titlesec}
\titleformat*{\section}{\large\bfseries}

\usepackage{fancyhdr}
\pagestyle{fancy}
\rhead{\small \color{gray} Notițe de Adrian Manea}
\lhead{\small \color{gray} IntProb-FIA, 2017---2018, L. Lemnete \& M. Olteanu}


\usepackage[autostyle = false, style = german]{csquotes}
\usepackage[romanian]{babel}
\newcommand\qq{\enquote}

%% COMENZILE MELE
\renewcommand*{\proofname}{Dem.:}
\newcommand\bb{\mathbb}
\newcommand\dr{\mathrm}
\newcommand\kal{\mathcal}
\newcommand\fk{\mathfrak}
\newcommand\op{\oplus}
\newcommand\ot{\otimes}
\newcommand\ds{\displaystyle}
\newcommand\id{\indent}
\newcommand\nid{\noindent}
\newcommand\seq{\subseteq}
\newcommand\sneq{\subsetneq}
\newcommand\speq{\supseteq}
\newcommand\spneq{\supsetneq}
\newcommand\rar{\rightarrow}
\newcommand\Rar{\Rightarrow}
\newcommand\xrar{\xrightarrow}
\newcommand\lar{\leftarrow}
\newcommand\Lar{\Leftarrow}
\newcommand\xlar{\xleftarrow}
\newcommand\lrar{\leftrightarrow}
\newcommand\Lrar{\Leftrightarrow}
\newcommand\ideal{\trianglelefteq}
\newcommand\ai{\text{ a.\^{i}. }}
\newcommand\sk{(\textit{Schi\c{t}\u{a}}) }
\newcommand\rhup{\rightharpoonup}
\newcommand\rhdn{\rightharpoondown}
\newcommand\lhup{\leftharpoonup}
\newcommand\lhdn{\leftharpoondown}
\newcommand\vid{\emptyset}
\newcommand\wh{\widehat}
\newcommand\ol{\overline}
\newcommand\NN{\mathbb{N}}
\newcommand\RR{\mathbb{R}}
\newcommand\ZZ{\mathbb{Z}}
\newcommand\QQ{\mathbb{Q}}
\newcommand\CC{\mathbb{C}}

%% STILURILE MELE
\newtheoremstyle{dotless}{}{}{\itshape}{}{\bfseries}{:}{ }{}

\theoremstyle{dotless}
\newtheorem{theorem}{Teorem\u{a}}[section]
\newtheorem{proposition}{Propozi\c{t}ie}[section]
\newtheorem{lemma}{Lem\u{a}}[section]
\newtheorem{corollary}{Corolar}[section]
\newtheorem{exercise}{Exerci\c{t}iu}

\newtheoremstyle{dotlessdr}{}{}{}{}{\bfseries}{:}{ }{}
\theoremstyle{dotlessdr}
\newtheorem{example}{Exemplu}[section]
\newtheorem{definition}{Defini\c{t}ie}[section]
\newtheorem{remark}{Observa\c{t}ie}[section]




\begin{document}

\begin{center}
  {\large\textbf{Seminar 7 \\
      Formula Green-Riemann}}
\end{center}

\vspace{1cm}

\section{Formula Green-Riemann}

Fie $ (K, \partial K) $ un compact cu bord orientat inclus în $ \RR^2 $ și
considerăm o 1-formă diferențială de clasă $ \kal{C}^1 $ pe o vecinătate a lui $ K $,
$ \alpha = P dx + Q dy $.
Atunci are loc \emph{formula Green-Riemann}:
\[
  \int_{\partial K} Pdx + Qdy =
  \iint_K \Big( \frac{\partial Q}{\partial x} - \frac{\partial P}{\partial y} \Big) dxdy.
\]

Această formulă ne permite să calculăm integrale curbilinii cu ajutorul
integralelor duble, însă doar în cazul în care forma diferențială are proprietățile
din ipoteză.

O consecință imediată este o formulă de calcul pentru arie:
\[
  \kal{A}(K) = \frac{1}{2} \int_{\partial K} xdy - ydx.
\]

%%%%%%%%%%%%%%%%%%%%%%%%%%%%%%%%%%%%%%%%%%%%%%%%%%%%%%%%%%%%%%%%%%%%%%

\section{Exerciții}

1. Calculați integralele duble:
\begin{enumerate}[(a)]
\item $ \ds\iint_D ydxdy $, unde $ D $ este mărginit de parabola $ y^2 = x $,
  cercul $ x^2 + y^2 - 2x = 0 $ și dreapta $ x = 2 $;
\item $ \ds\iint_D e^{x^2 + y^2} dxdy $, unde $ D = \{(x, y) \in \RR^2 \mid x^2 + y^2 \leq 1 \} $;
\item $ \ds\iint_D ydxdy $, unde $ D = \{(x, y) \in \RR^2 \mid (x - 2)^2 + y^2 \leq 1 \} $.
\end{enumerate}

\vspace{1cm}

2. Fie $ D \seq \RR^2 $ și $ f : D \to [0, \infty) $ o funcție continuă. Definim
mulțimea:
\[
  \Omega = \{ (x, y, z) \in \RR^3 \mid (x, y) \in D, 0 \leq z \leq f(x, y) \}.
\]
Să se calculeze volumul mulțimii $ \Omega $, folosind formula:
\[
  \kal{V}(\Omega) = \iint_D f(x, y) dxdy,
\]
în cazurile:
\begin{enumerate}[(a)]
\item $ D = \{(x, y) \in \RR^2 \mid x^2 + y^2 \leq 2y \}, f(x, y) = x^2 + y^2 $;
\item $ D = \{(x, y) \in \RR^2 \mid x^2 + y^2 \leq x, y > 0 \}, f(x, y) = xy $;
\item $ D = \{(x, y) \in \RR^2 \mid x^2 + y^2 \leq 2x + 2y - 1 \}, f(x, y) = y $.
\end{enumerate}

\vspace{1cm}

3. Calculați direct și aplicînd formula Green-Riemann integrala curbilinie
$ \ds\int_\Gamma \alpha $ în următoarele cazuri:
\begin{enumerate}[(a)]
\item $ \alpha = y^2 dx + xdy $, unde $ \Gamma $ este pătratul cu vîrfurile
  $ A(0, 0), B(2, 0), C(2, 2), D(0, 2)$;
\item $ \alpha = ydx + x^2 dy $, unde $ \Gamma $ este cercul cu centrul în
  origine și rază 2;
\item $ \alpha = ydx - xdy $, unde $ \Gamma $ este elipsa de semiaxe
  $ a $ și $ b $ și de centru $ O $.
\end{enumerate}

\vspace{1cm}

4. Fie forma diferențială:
\[
  \alpha = \frac{-y}{x^2 + y^2} dx + \frac{x}{x^2 + y^2} dy.
\]

Să se calculeze $ \int_\Gamma \alpha $, unde $ \Gamma $ este cercul cu centrul
în origine și rază 2.

\vspace{1cm}

5. Să se calculeze integrala curbilinie $ \ds\int_\Gamma xydx + \frac{x^2}{2} dy $,
pe conturul:
\[
  \Gamma = \{(x, y) \mid x^2 + y^2 = 1, x \leq 0 \leq y \} \cup %
  \{(x, y) \mid x + y = -1, x, y \leq 0 \}.
\]

\emph{Indicație:} curba $ \Gamma $ nu este închisă, deci nu putem aplica
formula Green-Riemann. Considerăm segmentul orientat $ [AB] $, cu $ A(0, -1) $ și
$ B(0, 1) $, cu care închidem curba, definind $ \Lambda = \Gamma \cup [AB] $.

Acum putem aplica formula Green-Riemann pe $ \Lambda $ și găsim:
\[
  \int_\Lambda xydx + \frac{x^2}{2} dy = \iint_K 0dxdy = 0,
\]
unde $ K $ este compactul mărginit de $ \Lambda $.

Atunci:
\[
  \int_\Gamma xydx + \frac{x^2}{2} dy = -\int_{[AB]} xydx + \frac{x^2}{2} dy = 0.
\]


\end{document}